\lesson{parallel}{多智能体：把独立工作拆开}
\diagram{p_parallel}
多智能体让多个 AI 协作。本章采用主智能体分工与汇总的方式：主会话统一目标，子智能体完成限定工作。子任务可分别阅读大量材料，再返回结论与证据入口，减少主会话中的杂乱输出。优先并行独立调查；依赖紧密的修改统一推进。
\topic{从两个只读助手开始}
\begin{enumerate}
\item \textbf{在主会话明确要求创建子智能体。}先确认工具支持，查看子任务运行记录。“只读”提示不是权限隔离：Codex 子任务沿用主会话的权限设置，不能靠角色说明单独限制；强制只读需相应设置主会话或独立运行环境。
\item \textbf{给各自不同任务。}一个查现有实现和可复用位置，另一个查相关方案、资料或验证风险。每个任务写清目标、材料位置、允许范围和返回内容，不假设它知道聊天历史或已读到项目规则。
\item \textbf{由主智能体统一判断。}依赖调查结论的修改，收齐依据后再做。并行写代码时约定文件范围、接口和完成条件，可用 Git worktree 分开工作目录；它不隔离权限或共享数据库。共享文件指定负责人，测试环境另行配置，合并后重查。
\end{enumerate}
\begin{prompt}{创建并行子任务}
目标：[当前任务]，共同材料：[需求、调研、计划与任务位置]。先核对工具能力和权限；支持时创建两个只读子任务：A查已有实现与复用位置，B查[方案、资料或验证风险]。材料不足先报缺口。\par
子任务不改文件、不扩展目标。若必须强制只读而无法配置，暂停启动并说明；角色提示不等于权限隔离。分别明确范围、输入、检查和返回格式：结论、位置、证据、未知项。主智能体在现有权限内推进独立工作；依赖结论的步骤等待结果，核对冲突后再推进。\par
若要并行修改，先明确授权、文件分工与合并检查。记录运行、完成或失败状态，未返回不算完成；不支持子智能体则返回可分别发送的任务说明，不模拟已启动。
\end{prompt}
分工让调查在各自上下文中进行，主会话保留决定与证据入口；它不能替代项目记录，也不保证消除信息损失。查看子任务是否实际启动，再收回来源、核对冲突。先从少量独立任务开始，比较总耗时、用量和返工后再扩量；智能体越多，重复工作与协调成本也可能越高。工具不支持时，可分会话调查后汇总。
